\documentclass[master, fontsize=12pt]{diploma}

\usepackage[reviewer]{review}

\student{Блинов Максим Алексеевич}
\faculty{№ 8 <<Компьютерные науки и прикладная математика>>}
\department{806}
\speciality{02.04.02 <<Фундаментальная информатика и информационные технологии>>}
\profile{Информатика}
\group{М8О-209СВ-24}

\theme{Оптимизация холодного старта и масштабирования моделей машинного обучения в бессерверном инференсе для обеспечения высокой доступности}
\workDate{\uline{\hspace{24pt}} июня \the\year\ года}

\supervisor{Булакина Мария Борисовна}
\supervisorCredentials{}

\firstConsultant{---}
\firstConsultantCredentials{}

\secondConsultant{---}
\secondConsultantCredentials{}

\reviewer{---}
\reviewerCredentials{}

\headOfDepartmentTitle{Заведующий кафедрой 806}
\headOfDepartment{Крылов Сергей Сергеевич}


\makeatletter
\renewcommand{\makeAffiliations}{
    \begingroup
    \fontsize{12pt}{14pt}\selectfont
    \setlength{\parskip}{0pt}
    \raggedright
    \noindent\textbf{Институт~(Филиал)} \expandafter\uline\expandafter{\expanded{\@faculty}\hfill}\par
    \noindent\textbf{Образовательный~центр~№} \uline{\@department\hfill}\par
    \noindent\textbf{Группа } \uline{\@group}\hspace{0.8em}\textbf{Направление~подготовки }%
    \expandafter\uline\expandafter{\expanded{\@speciality}\hfill}\par
    \noindent\textbf{\@profileType} \uline{\@profile\hfill}\par
    \noindent\textbf{Квалификация:} \uline{\textbf{\@degree}\hfill}\par
    \endgroup
}

\renewcommand{\makeReviewerLine}{
    \textbf{\@reviewerType} \uline{\@fullname{\@reviewType}, \@credentials{\@reviewType}\hfill}
}
\makeatother

\begin{document}
\fillReviewHeader
\setlength{\parskip}{0pt}
\sloppy

\textbf{Отмеченные достоинства:} Выпускная квалификационная работа посвящена актуальной задаче сокращения времени холодного старта и повышения доступности бессерверного инференса моделей машинного обучения в Kubernetes-среде. Актуальность выбранной темы обусловлена тем, что при использовании serverless-подхода значительная часть задержки первого запроса переносится из вычислительного контура в контур доставки данных, а рост объёма весов современных моделей делает повторную загрузку артефактов всё более затратной по времени и по внешнему трафику. Автор корректно формулирует исходное противоречие между требованиями к быстрому масштабированию и экономией ресурсов за счёт scale-to-zero. Сильной стороной работы является логично выстроенная последовательность инженерного анализа. Особо следует отметить практическую направленность диссертации. В работе не ограничиваются общими архитектурными рассуждениями: рассмотрены вопросы кэширования, репликации чанков, времени жизни метаданных и поведения системы при cache miss; проведено сопоставление предложенного подхода с emptyDir и NFS-сценариями по времени подготовки новых реплик, нагрузке на внешний источник и устойчивости к burst-нагрузке. Это позволяет сделать вывод о хорошей связи между теоретическим анализом, проектными решениями и их экспериментальной проверкой.
\\

\textbf{Отмеченные недостатки:} Экспериментальную часть целесообразно расширить за счёт большего числа моделей и более длительных нагрузочных сценариев. Дополнительного анализа также требует поведение системы при сетевых разбиениях и дальнейшем росте числа узлов кластера.
\\

\textbf{Заключение}: Работа выполнена на высоком научно-техническом уровне, сочетает анализ предметной области, проектирование распределённой архитектуры и обоснование её эффективности. Отмеченные недостатки носят рекомендательный характер. Работа удовлетворяет требованиям, предъявляемым к выпускным квалификационным работам по направлению «Фундаментальная информатика и информационные технологии», а автор Блинов М. А. заслуживает оценки «отлично»  и присвоения квалификации инженер-исследователь.

\makeReviewSignature
\end{document}
